	\section{Pseudorandom generators}\label{sec:prgs}
	
	Some of our PRGs are obtained from a general framework that allows us to reduce the task of fooling $\FORMULA\circ\mathcal{G}$ to the task of fooling the class of functions which are the parity or conjunction of few functions from $\mathcal{G}$.
	
	\subsection{The general framework}
	
	We show that in order to get a PRG for the class of subquadratic-size formulas with leaf gates in $\mathcal{G}$, it suffices to get a PRG for very simple sublinear-size formulas: either $\XOR\circ\mathcal{G}$ or $\AND\circ\mathcal{G}$.  
	
	\begin{theorem}[PRG for $\FORMULA\circ\mathcal{G}$ from PRG for $\XOR\circ\mathcal{G}$ or $\AND\circ\mathcal{G}$]\label{thm:prg-main}
		Let $\mathcal{G}$ be a class of gates on $n$ bits. For any integer $s>0$ and any $0<\varepsilon<1$, there exists a constant $c>0$ such that the following holds. If a distribution $\mathcal{D}$ over $\boolpm^n$ $\left(2^{-c\cdot\sqrt{s}\cdot\log (s)\cdot \log(1/\varepsilon)}\right)$-fools the $\XOR$ (parity) or the $\AND$ (conjunction) of $c\cdot\sqrt{s}\cdot \log(1/\varepsilon)$ arbitrary functions from $\mathcal{G}$, then $\mathcal{D}$ also $\varepsilon$-fools $\FORMULA[s]\circ\mathcal{G}$.
	\end{theorem}
	
	\begin{proof}\renewcommand{\qedsymbol}{$\square$ (\Cref{thm:prg-main})}
    	We first show the case where $\mathcal{D}$ fools the parity of a few functions from $\mathcal{G}$. The proof can be easily adapted to the case of conjunction.
    	
    	Let $C=F(g_1,g_2\dots,g_s)$ be a function in $\FORMULA[s]\circ\mathcal{G}$, where $F$ is a formula, and $g_1,g_2,\dots,g_s$ are functions from the class $\mathcal{G}$. Let $\uniform$ be the uniform distribution over $\boolpm^n$. 
    	We need to show
    	\begin{equation}\label{eq:prg-1}
    		\Exp[C(\mathcal{D})] \myapprox{\varepsilon} \Exp[C(\uniform)].
    	\end{equation}
    	Let $p$ be a $(\varepsilon/3)$-approximating polynomial for $F$ given by \Cref{thm:approx-polynomial}. Note that the degree of $p$ is
    	\[
    		d=O(\sqrt{s}\cdot\log(1/\varepsilon)).
    	\] 
    	Let us replace $F$, the formula part of $C$, with $p$ and let
    	\[
    		\tilde{C}\vcentcolon=p(g_1,g_2\dots,g_s).
    	\]
    	Since $\tilde{C}$ point-wisely approximates $C$, we have
    	\[
    		\Exp[\tilde{C}(\uniform)] \myapprox{\varepsilon/3} \Exp[C(\uniform)],
    	\]
    	and
    	\[
    		\Exp[\tilde{C}(\mathcal{D})] \myapprox{\varepsilon/3} \Exp[C(\mathcal{D})].
    	\]
    	Then to show \Cref{eq:prg-1}, it suffices to show
    	\[
    		\Exp[\tilde{C}(\mathcal{D})] \myapprox{\varepsilon/3} \Exp[\tilde{C}(\uniform)].
    	\]
    	We have
    	\begin{align*}
        	\label{eq:prg-2}
        	\Exp_{x\sim \mathcal{D}}[\tilde{C}(x)]&=\Exp_{x\sim D}\left[\sum_{\substack{S\subseteq[s]: \\ |S|\leq d}}\hat{p}(S)\cdot\prod_{i\in S}g_i(x)\right]\\
        	&=\sum_{\substack{S\subseteq[s]: \\ |S|\leq d}}\hat{p}(S)\cdot\Exp_{x\sim \mathcal{D}}\left[\prod_{i\in S}g_i(x)\right].
        	\numberthis
    	\end{align*}
    	Now note that for each $S\subseteq[s]$, $\prod_{i\in S}g_i(x)$ computes the $\XOR$ of at most $d$ functions from $\mathcal{G}$.
    	Using the fact the distribution $\mathcal{D}$ $\left(\delta=1/2^{c\cdot\sqrt{s}\cdot\log (s)\cdot \log(1/\varepsilon)}\right)$-fools the $\XOR$ of any $d$ functions from $\mathcal{G}$, we get
    	\begin{align*}
        	\Exp_{x\sim \mathcal{D}}[\tilde{C}(x)]&=\sum_{\substack{S\subseteq[s]: \\ |S|\leq d}}\hat{p}(S)\cdot\Exp_{x\sim D}\left[\prod_{i\in S}g_i(x)\right]\\
        	&=\sum_{\substack{S\subseteq[s]: \\ |S|\leq d}}\hat{p}(S)\cdot\left(\Exp_{x\sim \uniform}\left[\prod_{i\in S}g_i(x)\right]+\delta_S\right) \tag{where $|\delta_S|\leq \delta$} \\
        	&=\sum_{\substack{S\subseteq[s]: \\ |S|\leq d}}\left(\hat{p}(S)\cdot\Exp_{x\sim \uniform}\left[\prod_{i\in S}g_i(x)\right]+\hat{p}(S)\cdot\delta_S\right)\\
        	&=\sum_{\substack{S\subseteq[s]: \\ |S|\leq d}}\hat{p}(S)\cdot\Exp_{x\sim \uniform}\left[\prod_{i\in S}g_i(x)\right]+\sum_{\substack{S\subseteq[s]: \\ |S|\leq d}} \hat{p}(S)\cdot\delta_S\\
        	&=\Exp_{x\sim \uniform}[\tilde{C}(x)]+\sum_{\substack{S\subseteq[s]: \\ |S|\leq d}} \hat{p}(S)\cdot \delta_S.
    	\end{align*}
    	It remains to show
    	\[
    		\left| \sum_{\substack{S\subseteq[s]: \\ |S|\leq d}} \hat{p}(S)\cdot \delta_S \right| \leq \varepsilon/3.
    	\]
    	Note that because $p(z)\in[1-\varepsilon/3,1+\varepsilon/3]$ for every $z\in\boolpm^s$, we have
    	\[
    		|\hat{p}(S)|=\left|\Exp_{z\sim \boolpm^s}\left[p(z)\cdot\prod_{i\in S} z_i\right]\right| \leq 1+\varepsilon/3 < 2.
    	\]
    	Then,
    	\[
    		\left|\sum_{\substack{S\subseteq[s]: \\ |S|\leq d}} \hat{p}(S)\cdot \delta_S \right| \leq  \sum_{\substack{S\subseteq[s]: \\ |S|\leq d}} |\hat{p}(S)|\cdot |\delta_S|\leq \delta \cdot \sum_{\substack{S\subseteq[s]: \\ |S|\leq d}} |\hat{p}(S)| \leq \delta\cdot s^{O(\sqrt{s}\cdot\log(1/\varepsilon))} \leq \varepsilon/3,
    	\]
    	where the last inequality holds for some sufficiently large constant $c$.
    	
    	To show the case of conjunction, we can write the approximating polynomial as the sum of all degree-$d$ monomials, each of which is the $\AND$ of at most $d$ variables. One way to do this is to use the domain $\bool$ instead of $\boolpm$ in the above argument. We need to show that the coefficients in this case still have small magnitude.
    	
    	\begin{claim}\label{claim:prg-magnitude}
    		Let $p\colon\boolpm^{n}\to\real$ be a degree-$d$ polynomial of the form
    		\[
    			p(x)=\sum_{\substack{S\subseteq[n]: \\ |S|\leq d}} \hat{p}(S)\cdot \prod_{i\in S} x_i,
    		\]
    		and let $q\colon\bool^{n}\to\real$ be the corresponding polynomial of $p$ over the domain $\{0,1\}^n$, of the form
    		\[
    			q(y)=\sum_{\substack{T\subseteq[n]: \\ |T|\leq d}} \hat{q}(T)\cdot \prod_{i\in T} y_i.
    		\]
    		Then,
    		\[
    			|q|_1 = \sum_{\substack{T\subseteq[n]: \\ |T|\leq d}} |\hat{q}(T)| \leq n^{O(d)}\cdot \max_{\substack{S\subseteq[n]: \\ |S|\leq d}} |\hat{p}(S)|.
    		\]
    	\end{claim}
    	
    	\begin{proof}\renewcommand{\qedsymbol}{$\square$ (\Cref{claim:prg-magnitude})}
    		We have
    		\begin{align*}
    			q(y_1,y_2,\dots,y_n) &= p(1-2y_1,1-2y_2,\dots,1-2y_n)\\
    			&= \sum_{\substack{S\subseteq[n]: \\ |S|\leq d}} \hat{p}(S)\cdot \prod_{i\in S} (1-2y_i)\\
    			&= \sum_{\substack{S\subseteq[n]: \\ |S|\leq d}} \hat{p}(S)\cdot \left(\sum_{\ell\in \{0,1\}^{|S|}}\prod_{\substack{j\in S:\\ \ell_j=1}}-2y_j \right)\\
    			&= \sum_{\substack{S\subseteq[n]: \\ |S|\leq d}} \sum_{\ell\in \{0,1\}^{|S|}} \hat{p}(S)\cdot(-2)^{|\ell|}\cdot \prod_{\substack{j\in S:\\ \ell_j=1}} y_j. \tag{where $|\ell| = \sum_{i=1}^{|S|}\ell_i$}
    		\end{align*}
    		For a pair $(S,\ell)$ where $S \subseteq[n]$, $|S|\leq d$ and $\ell\in\{0,1\}^{|S|}$, let us define the polynomial $q_{(S,\ell)}$ as
    		\[
    			q_{(S,\ell)}(y)= \hat{p}(S)\cdot(-2)^{|\ell|}\cdot \prod_{\substack{j\in S:\\ \ell_j=1}} y_j.
    		\]
    		Note that there are at most $n^{d}\cdot2^{d}$ many pairs of such $(S,\ell)$'s and for each $(S,\ell)$, we have
    		\begin{align*}
    			|q_{(S,\ell)}|_1 = \left|\hat{p}(S)\cdot(-2)^{|\ell|}\right|\leq 2^{d}\cdot |\hat{p}(S)|.
    		\end{align*}
    	    Finally we have
    		\[
    			|q|_1 = \left|\sum_{(S,\ell)} q_{(S,\ell)}\right|_1\leq \sum_{(S,\ell)} |q_{(S,\ell)}|_1\leq n^{d}\cdot 2^d \cdot 2^d \cdot \max_{\substack{S\subseteq[n]: \\ |S|\leq d}} |\hat{p}(S)|,
    		\]
    		as desired.
    	\end{proof}
    	
    	This completes the proof of \Cref{thm:prg-main}.
	\end{proof}

    \subsection{Formulas of low-communication functions in the number-in-hand setting}
    
    In this subsection, we will use $\bool$ as the Boolean basis.
    
	\begin{theorem}\label{thm:prg-formula-lc}
		For any integers $k\geq 2$, $s>0$ and any $0<\varepsilon<1$, let $\mathcal{G}$ be the class of functions that have $k$-party number-in-hand $(\varepsilon/6s)$-error randomized communication protocols of cost at most $R$. There exists a PRG that $\varepsilon$-fools $\FORMULA[s]\circ \mathcal{G}$ with seed length
		\[
			n/k+O\left(\sqrt{s}\cdot (R +\log (s))\cdot \log(1/\varepsilon) +\log(k)\right)\cdot\log(k).
		\]
	\end{theorem}
    
    We need the following PRG that fools single functions with low communication complexity in the number-in-hand model. The proof is presented in \Cref{appendix:prg_cc} (\Cref{lem:prg-lc-app}) for completeness.

	\begin{theorem}[\cite{ASWZ96,INW94}]\label{thm:prg-lc}
		For any $k\geq 2$, there exists a PRG that $\delta$-fools any $n$-bits functions with $k$-party number-in-hand deterministic communication complexity of at most $D'$, with seed length
		\[
			n/k+O\left(D' + \log(1/\delta)+\log(k)\right)\cdot\log(k).
		\]
	\end{theorem}

    Next, we show a PRG for $\FORMULA\circ\mathcal{G}$, where $\mathcal{G}$ is the class of functions with low-cost communication protocols in the number-in-hand setting. We first show for the case of deterministic protocols.
    
    \begin{theorem}\label{thm:prg-formula-lc-det}
		For any integers $k\geq 2$ and $s>0$, let $\mathcal{G}$ be the class of functions whose $k$-party number-in-hand deterministic communication complexity are at most $D$. There is a PRG that $\varepsilon$-fools $\FORMULA[s]\circ \mathcal{G}$ of size $s$ with seed length
		\[
			n/k+O\left(\sqrt{s}\cdot \log(1/\varepsilon)\cdot (D +\log (s)) +\log(k)\right)\cdot\log(k).
		\]
	\end{theorem}
	
	\begin{proof}
		By \Cref{thm:prg-main}, it suffices to show a PRG that $\left(\delta=1/2^{c\cdot\sqrt{s}\cdot\log (s)\cdot \log(1/\varepsilon)}\right)$-fools every function that is the $\XOR$ of $t=c\cdot\sqrt{s}\cdot \log(1/\varepsilon)$ arbitrary functions from $\mathcal{G}$. Note that such a function has deterministic communication complexity at most $D'=t\cdot D$. Then \Cref{thm:prg-formula-lc-det} follows from \Cref{thm:prg-lc}.
	\end{proof}
	
	We now establish the randomized case.
    
    \begin{proof}[Proof of \Cref{thm:prg-formula-lc}]
    	Let $C$ be a function in $\FORMULA[s]\circ \mathcal{G}$. For each of the leaf functions in $C$, consider a $k$-party number-in-hand randomized protocol of cost at most $R$ that has an error at most $\varepsilon/(6s)$. By taking a union bound over the $s$ leaf functions and by viewing a randomized protocol as a distribution of deterministic protocols (as shown in the proof of \Cref{claim:lb-approximating}), we get the following which is a (point-wisely) $(\varepsilon/3)$-approximating function for $C$:
        \[
            \tilde{C}(x) \vcentcolon=\sum_{i} p_i \cdot D_i(x),
        \]
        where each $p_i\in[0,1]$ is some probability density value (so $\sum_i p_i =1$), and each $D_i$ is a formula whose leaves are functions with \emph{deterministic} communication complexity at most $R$. Then to $\varepsilon$-fool $C$, it suffices to $(\varepsilon/3)$-fool its $(\varepsilon/3)$-approximating function $\tilde{C}$. Also, since $\tilde{C}$ is a convex combination of the $D_i$'s, it suffices to $(\varepsilon/3)$-fools all the $D_i$'s. We will do this using the PRG form \Cref{thm:prg-formula-lc-det}. We get that there exists a PRG that $(\varepsilon/3)$-fools each $D_i$ with seed length
        \[
            n/k+O\left(\sqrt{s}\cdot(R+\log(s))\cdot\log(1/\varepsilon)+\log(k)\right)\cdot\log(k),
        \]
        as desired.
    \end{proof}

%
   
%
   
%
%
%
%
%
%
	
%
	
%
	
%
%
%
%
%
%
	
%
	
%
%
%

%
%
%
%
%
%
%
%
%
%
%
%
%
%
%
%
%
%
%
%
%
%
%
%
%
%
%
%
%
%
%
%
%
%
%
%
%
%
%
%
%
%
%
%
%
%
%
%
%
%
%
%
%
    
%
	
%
    
%
%
%
%
%
%
	
%
%
%
	
%
    
%
%
%
%
%
%
%
%
%
%
%

    \subsection{Applications: Fooling formulas of SYMs, LTFs, XORs, and AC\texorpdfstring{$^0$ circuits}{AC0}}

    \subsubsection{\texorpdfstring{$\FORMULA\circ\SYM$}{FORMULA-SYM} and \texorpdfstring{$\FORMULA\circ\LTF$}{FORMULA-LTF}}
	
	Here, we show how the PRG in \Cref{thm:prg-formula-lc} implies PRGs for $\FORMULA\circ\LTF$ and $\FORMULA\circ\SYM$.
	
	\begin{theorem}\label{thm:prg-formula-sym}
		For any size $s>0$ and $0<\varepsilon<1$, there exists a PRG that $\varepsilon$-fools $\FORMULA[s]\circ \LTF$ with seed length
        \[
			O\left(n^{1/2}\cdot s^{1/4}\cdot \log(n)\cdot \log(n/\varepsilon)\right).
		\]
		For $\FORMULA[s]\circ \SYM$, the seed length is
		\[
		    O\left(n^{1/2}\cdot s^{1/4}\cdot \log(n)\cdot \log(1/\varepsilon) \right).
		\]
	\end{theorem}
	
	We need the fact that the class of $\LTF$ has low communication complexity in the number-in-hand model. Consider the following $k$-party $\SG_m$ problem where the $i$-th party holds a $m$-bit number $z_i$ in hand and they want to determine whether $\sum_{i=1}^k z_i > \theta$, where $\theta$ is a fixed number known to all the parties. Nisan~\cite{Nis94} gave an efficient randomized protocol (with public randomness) for this problem.
    
%
%
%
    
    \begin{theorem}[\cite{Nis94}\footnote{Viola~\cite{Vio15} gave a $\delta$-error randomized protocol for the $k$-party $\SG_m$ problem of cost $O(k\cdot\log(k)\cdot\log(m/\delta))$, which is better than Nisan's protocol when $k=m^{o(1)}$.}]\label{thm:prg-k-party-sum}
	    Let $m>0$ be an integer. For any integer $2\leq k\leq m^{O(1)}$, and any $0<\delta<1$, there exists a $\delta$-error randomized protocol of cost $O(k\cdot\log(m)\cdot\log(m/\delta))$ for the $k$-party $\SG_m$ problem.
	\end{theorem}
	    
    By \Cref{thm:prg-k-party-sum} and the fact that every linear threshold function on $n$ bits has a representation such that the weights are $O(n\log(n))$ integers~\cite{MTT61}, we get the following.
    
    \begin{corollary}\label{cor:prg-k-party-LTF}
         For every $k\geq 2$ and $0<\delta<1$, the $k$-party number-in-hand $\delta$-error randomized communication complexity of $\LTF$ is $O(k\cdot\log(n)\cdot\log(n/\delta))$.
    \end{corollary}
    
	\begin{proof}[Proof of \Cref{thm:prg-formula-sym}]
	    By \Cref{cor:prg-k-party-LTF} and \Cref{thm:prg-formula-lc}, for \emph{every} $k\geq 2$ we get a PRG for $\FORMULA\circ\LTF$ of seed length
        \[
    		n/k+O\left(\sqrt{s}\cdot k\cdot \log(n) \cdot\log(ns/\varepsilon)\cdot\log(1/\varepsilon)+\log(k)\right)\cdot\log(k).
    	\]
	    By choosing 
	    \[
	        k=\frac{n^{1/2}}{s^{1/4}\cdot\log(n)\cdot\log(n/\varepsilon)},
	    \]
	    the claimed seed length follows from a simple calculation.
	
	    For $\FORMULA\circ\SYM$, note that every $n$-bit symmetric function has a deterministic $k$-party number-in-hand communication protocol of cost at most $k\cdot\log(n)$. Then the rest can be shown using a similar argument as above (by choosing $k=n^{1/2}/\left(s^{1/4}\cdot\log(n)\right)$).
	\end{proof}

%
       
%

%
%
%
%
%
%

    \subsubsection{\texorpdfstring{$\FORMULA\circ\XOR$}{FORMULA-XOR}}
    
    For the case of $\FORMULA\circ\XOR$, we get a PRG with better seed length. 
    
	\begin{theorem}\label{thm:prg-formula-xor}
		For any size $s>0$ and $0<\varepsilon<1$, there exists a PRG that $\varepsilon$-fools $\FORMULA[s]\circ \XOR$ with seed length
		\[
			O\left(\sqrt{s}\cdot\log(s)\cdot\log(1/\varepsilon)+\log(n)\right).
		\]
	\end{theorem}
	
	\begin{proof}
	    By \Cref{thm:prg-main}, to fool $\FORMULA[s]\circ\mathcal{G}$, it suffices to $\left(\delta=1/2^{O\left(\sqrt{s}\cdot\log (s)\cdot \log(1/\varepsilon)\right)}\right)$-fool the $\XOR$ of a few functions from $\mathcal{G}$, where $\mathcal{G}$ in this case is the set of all $\XOR$ functions. Note that the $\XOR$ of any set of $\XOR$ functions simply computes some $\XOR$ function. Therefore, we can use small-bias distribution, which fools every $\XOR$ function, to fool $\FORMULA[s]\circ\XOR$. Finally, note that there are known constructions for $\delta$-bias distributions that use $O(\log(n/\delta))$ random bits (see e.g.~\cite{AGHP92}).
	\end{proof}
	
%
%
%
%
%
%
%
%
%
%
    \subsubsection{\texorpdfstring{$\FORMULA\circ\AC^0$}{FORMULA-AC0}}

	Another application of \Cref{thm:prg-main} is to take $\mathcal{G}$ to be the set all functions that can be computed by small constant-depth circuits ($\AC^0$). Note the state-of-the-art PRG against size-$M$ depth-$d$ $\AC^0$ has a seed length of $\log^{d+O(1)}(Mn)\cdot\log(1/\varepsilon)$~\cite{ST19}. Below, let $\AC^0_{d,M}$ denote the class of depth-$d$ circuits of size at most $M$.
	
	\setcounter{theorem}{34}
\begin{theorem}\label{thm:prg-formula-ac0}
		For any size $s,m>0$ and $0<\varepsilon<1$, there exists a PRG that $\varepsilon$-fools $\FORMULA\circ \AC^0_{d,M}$ of size $s$ with seed length
		\[
			\log^{d+O(1)}(Mn)\cdot\sqrt{s}\cdot\log(s)\cdot\log(1/\varepsilon).
		\]
	\end{theorem}
	
%
%
%
%
%
%
%
%
%
	\subsection{Formulas of low number-on-forehead communication leaf gates}
	
	In this section, we show a PRG with mild seed length for formulas of functions with low \emph{multi-party number-on-forehead} communication complexity.
	
	\setcounter{theorem}{36}
\begin{theorem}\label{thm:mild-prg}
	    Let $\mathcal{G}$ be a class of $n$-bits functions. For any size $s>0$, there exists a PRG that $\varepsilon$-fools $\FORMULA[s]\circ \mathcal{G}$, with seed length
		\[
		    n-\frac{n}{O\left(\sqrt{s}\cdot k\cdot4^k\cdot \left( R^{(k)}_{\varepsilon/(2s)}(\mathcal{G})+\log(n)\right)\cdot\log(n/\varepsilon)\right)}.
		\]
	\end{theorem}
	
	The PRG is constructed using the hardness vs. randomness paradigm.
	    
	\subsubsection{Hardness based PRGs}
	
	We show how to construct the PRG using the average-case hardness result for formulas of functions with low multi-party communication complexity (\Cref{thm:lb-main}). We start with some notations. For $x\in\boolpm^m$ and an integer $k$ such that $k$ divides $m$, we consider a partition of $x$ into $k$ equal-sized consecutive blocks and write $x=x^{(1)},x^{(2)},\dots,x^{(k)}$, where $x^{(i)}\in\boolpm^{m/k}$ for each $i\in [k]$.
    	
	\begin{lemma}\label{lem:mild-prg}
	    For any integers $m,t,k>0$ such that $k$ divides $m,t$, let $\mathcal{G}$ be a class of functions on $mt+t$ bits, and let $G\colon\boolpm^{m\times t}\to\boolpm^{mt+t}$ be
		\begin{align*}
			&\quad G(x_1,x_2,\dots,x_t) \\
			&=\left(x_1^{(i)},x_2^{(i)},\dots,x_t^{(i)},\GIP_{m}^{k}\left(x_{(i-1)\cdot(t/k)+1}\right),\GIP_{m}^{k}\left(x_{(i-1)\cdot(t/k)+2}\right),\dots, \GIP_{m}^{k}\left(x_{i\cdot(t/k)+1}\right)\right)_{i\in[k]},
		\end{align*}
		where $x_1,x_2,\dots,x_t\in\boolpm^m$.
		Then $G$ is a PRG that $(t\cdot \varepsilon)$-fools $\FORMULA \circ \mathcal{G}$ of size 
		\[
			s=\Omega\left( \frac{m^2}{k^2\cdot16^{k}\cdot\left(R^{(k)}_{\varepsilon/(2m^2)}(\mathcal{G})+\log m\right)^2\cdot\log^2(1/\varepsilon)}\right).
		\]
	\end{lemma}
	
	\begin{proof}
		The high level idea is as follows. We argue that if there is a $\FORMULA \circ \mathcal{G}$ of the claimed size that breaks the PRG, then there is a $\FORMULA \circ \mathcal{G'}$ of the same size that computes $\GIP$ on $m$ bits, where $\mathcal{G'}$ has a $k$-party communication complexity that is at most that of $\mathcal{G}$ \emph{with respect to the $m$-bit input}, and hence contradicts the $\FORMULA \circ \mathcal{G'}$ complexity of the generalized inner product function. The resulting formula is obtained by fixing some input bits of the original $\FORMULA \circ \mathcal{G}$ which breaks the PRG.

		We use a hybrid argument. First consider the distribution given by $G$, where we replace each $\GIP(x_j)$ ($j\in[t]$) with a uniformly random bit; let us denote those random bits as $\uniform_{j}$ for $j\in[t]$ (note that this is just the uniform distribution). Then for each $j\in[t]$, define $H_j$ to be the distribution that we substitute back $\GIP(x_1),\GIP(x_2),\dots,\GIP(x_{j})$ for the corresponding uniform bits in the previous distribution. 

		For the sake of contradiction, suppose there exists a $\FORMULA \circ \mathcal{G}$ $C$ of size $s$ such that
		\[
			\left| \Prob[C(H_t)=1]-\Prob[F(H_0)=1] \right| > t\cdot\varepsilon.
		\]
		By the triangle inequality, there exists a $1\leq j\leq k$ such that
		\[
			\left| \Prob[C(H_j)=1]-\Prob[C(H_{j-1})=1] \right| >\varepsilon.
		\]
		Then by averaging, there exist some fixings of $x_1,\dots,x_{j-1},x_{j+1},\dots,x_{t}$ and $\uniform_{j+1},\dots,\uniform_{t}$ to $C$ such that the above inequality still holds. Let us denote by $C'$ the circuit obtained by $C$ after such fixings and assume without loss of generality $(k-1)t/k\leq j \leq t$. Then we have
		\begin{equation}\label{eq:distinguish}
		    \left| \Prob\left[C'\left(x_j^{(1)},x_j^{(2)},\dots,x_j^{(k)},\GIP(x_j)\right)=1\right]-\Prob\left[C'\left(x_j^{(1)},x_j^{(2)},\dots,x_j^{(k)},\uniform_j\right)=1\right] \right| > \varepsilon.
		\end{equation}
		By a standard ``unpredictability implies pseudorandomness" argument~\cite{Yao82}, we can show that there is some circuit $C''$, obtained from $C'$ by fixing some value for the last bit, that computes the generalized inner product function on $m$ bits with probability greater than $1/2+\varepsilon$ over uniformly random inputs. Note that the size of $C''$ is the same as $C'$ (hence also $C$) , and also $C''$ can be computed by some $\FORMULA \circ \mathcal{G'}$, where $R^{(k)}_{\delta}(\mathcal{G'})\leq R^{(k)}_{\delta}(\mathcal{G})$ for every $\delta$. This contradicts hardness of $\GIP$ for such circuits (\Cref{thm:lb-main}).
	\end{proof}
	
	We are now ready to prove \Cref{thm:mild-prg}.
	
	\begin{proof}[Proof of \Cref{thm:mild-prg}]
		Consider \Cref{lem:mild-prg}. Let $n=mt+t$, and we have $m=\left(\frac{n}{t}-1\right)$. Then 
		\Cref{lem:mild-prg} gives a PRG that $\varepsilon$-fools $\FORMULA\circ\mathcal{G}$ of size
		\begin{align*}
			s &= \Omega\left( \frac{m^2}{k^2\cdot16^{k}\cdot\left(R^{(k)}_{\varepsilon/(2m^2)}(\mathcal{G})+\log m\right)^2\cdot\log^2(t/\varepsilon)}\right)\\
			&\geq \Omega\left(\left(\frac{n}{t}\right)^2/\left(k^2\cdot16^{k}\cdot\left(R^{(k)}_{\varepsilon/(2n^2)}(\mathcal{G})+\log n\right)^2\cdot\log^2(n/\varepsilon)\right)\right),
		\end{align*}
		which yields
		\[
			t\geq \Omega\left(\frac{n}{\sqrt{s}\cdot k\cdot4^{k}\cdot\left(R^{(k)}_{\varepsilon/(2n^2)}(\mathcal{G})+\log n\right)\cdot\log(n/\varepsilon)}\right).
		\]
		Note that the seed length in this case is $n-t$.
	\end{proof}

	\subsubsection{\texorpdfstring{$\MKtP$}{MKtP} lower bounds}
	
	The PRG in \Cref{thm:mild-prg} is sufficient to give an $\MKtP$ lower bound for formulas of functions with low multi-party communication complexity.
	
	\begin{theorem}\label{thm:prg-mktp}
	    For any integer $s>0$ and any class of $N$-bit function $\mathcal{G}$, if $\MKtP$ on $N$-bit can be computed by some function $\FORMULA[s]\circ \mathcal{G}$, then
	    \[
            s= \frac{N^2}{k^2\cdot16^{k}\cdot R^{(k)}_{1/3}(\mathcal{G})\cdot \polylog(N)}.
	    \]
	\end{theorem}
	
	\begin{proof}
    	Let $C$ be a function in $\FORMULA\circ\mathcal{G}$ of size less than
    	\[
    	    \frac{N^2}{k^2\cdot16^{k}\cdot R^{(k)}_{1/3}(\mathcal{G})\cdot\log^c(N)}
    	\]
    	where $c>0$ is some sufficiently large constant. By \Cref{thm:mild-prg}, we have that there is a PRG that $(1/3)$-fools $C$ and its seed length is
    	\[
             N-\polylog(N).
    	\]
    	Also, since the PRG is polynomial-time computable, we get that for every seed, the output of the PRG has $\mathrm{Kt}$ complexity at most $\theta=N-\polylog(N)$. However, consider the $\MKtP$ function with a threshold parameter $\theta$; this function is not fooled by such a PRG, since it accepts every output of the PRG and rejects a uniformly random string with high probability. 
    \end{proof}
    
